\documentclass[11pt,a4paper]{article}
\usepackage[T1]{fontenc}
\usepackage{lmodern,amsmath,amssymb,amsthm,mathtools}
\usepackage[margin=25mm]{geometry}
\usepackage{booktabs,longtable,array,microtype}
\usepackage{xcolor,fancyhdr,enumitem}
\usepackage[unicode,colorlinks=true,linkcolor=blue!45!black,urlcolor=blue!45!black,citecolor=blue!45!black]{hyperref}
\usepackage{xurl}
\setcounter{MaxMatrixCols}{20}
\newtheorem{theorem}{Theorem}[section]
\newtheorem{lemma}[theorem]{Lemma}
\newtheorem{proposition}[theorem]{Proposition}
\newtheorem{corollary}[theorem]{Corollary}
\theoremstyle{definition}\newtheorem{definition}[theorem]{Definition}
\theoremstyle{remark}\newtheorem{remark}[theorem]{Remark}
\newcommand{\Q}{\mathbb Q}\newcommand{\Z}{\mathbb Z}\newcommand{\F}{\mathbb F}
\newcommand{\OO}{\mathcal O}\newcommand{\Cl}{\operatorname{Cl}}
\newcommand{\Sel}{\operatorname{Sel}}\newcommand{\rank}{\operatorname{rank}}
\newcommand{\Sha}{\mathop{\mbox{$\amalg$}}}\newcommand{\Nm}{\operatorname{N}}
\newcommand{\file}[1]{\texttt{\nolinkurl{#1}}}
\setlength{\emergencystretch}{3em}
\setlist{itemsep=3pt,topsep=5pt}


\title{Primitive solutions of a generalized Fermat equation of signature (17, 17, 3)}
\author{Yuwen Wang\texorpdfstring{\textsuperscript{a,*}}{}}
\date{}
\begin{document}
\maketitle
\begin{center}
\textsuperscript{a}School of Mechanical Engineering, Xi'an Jiaotong University\\
No. 28 Xianning West Road, Xi'an, Shaanxi 710049, P. R. China\\
* Corresponding author: Yuwen Wang. E-mail: \texttt{3010317@proton.me}\\
ORCID: 0009-0000-3188-9605
\end{center}



\begin{abstract}
We determine the primitive nonzero integer solutions of $X^3=Y^{17}+2Z^{17}$ and show that they are exactly $(1,-1,1)$ and $(-1,1,-1)$. The proof combines the modular method with an explicit classification of points on $X_0(17)$ over the degree-$18$ line-stabilizer field of a fixed mod-$17$ Galois representation. An unconditional ideal-class computation supplies the complete squareclass space for isogeny descent. Lower and higher Cassels--Tate pairings, together with three independent points, determine the Mordell--Weil rank to be $3$. Saturation identifies the required $5$-adic closure. Exact calculations modulo $25$ and integral analytic uniqueness on the remaining residue disks determine all modular points satisfying the prescribed rationality and congruence conditions. Their modular invariants form a set of three rational numbers. Level lowering and local comparison force the invariant of every putative nontrivial primitive solution to belong to this set. Three polynomial factorizations and elementary valuation arguments then complete the classification. As consequences, we obtain several primitive nonexistence results and show that the affine curve $x^3-y^{17}=2$ has the unique rational point $(1,-1)$. The arithmetic data and verification programs accompany the paper.
\end{abstract}

\noindent\textbf{Keywords and phrases:} generalized Fermat equations, modular curves, isogeny descent, Cassels--Tate pairing, elliptic Chabauty, rational points.

\noindent\textbf{2020 Mathematics Subject Classification:} 11D41, 11G05, 11G18, 11R29, 11Y50.


\section{Introduction}

We determine all primitive nonzero integer solutions of the generalized Fermat equation of signature $(17,17,3)$. Our main result is the following theorem, proved in Section~\ref{sec:applications} from the modular-point classification of Proposition~\ref{prop:rationalj}.

\begin{theorem}[Primitive ternary classification]\label{thm:ternary}
The only nonzero integer solutions of $X^3=Y^{17}+2Z^{17}$ with $\gcd(Y,Z)=1$ are $(1,-1,1)$ and $(-1,1,-1)$. In particular there is no primitive positive integer solution.
\end{theorem}



Two completeness statements connect the global arithmetic to the local calculation. The unconditional class-group bound makes the squareclass search exhaustive. Saturation identifies the closure of the subgroup generated by the explicit points and rational torsion with that of the full Mordell--Weil group. The calculation modulo $25$ determines the possible local parameters, and integral analytic uniqueness resolves the entire $5$-adic residue disks.

The arithmetic contribution is the explicit determination, under the chosen completion at five, of the modular points over the degree-18 line-stabilizer field whose invariants lie in the rational 5-adic field and are congruent to two modulo five. The class-group calculation, complete local Kummer images and two descent pairings give the rank computation needed for the local Mordell--Weil closure. Determining these points links the residual representation to the final Diophantine contradiction.

Generalized Fermat equations connect perfect-power conditions with the arithmetic of curves and their Galois representations. Kraus gives a survey of this setting \cite{Kraus}. Related asymptotic questions over totally real fields are studied by Jha and Sahoo \cite{JhaSahoo}.

The modularity input has its origins in the work of Wiles and Taylor--Wiles on semistable elliptic curves \cite{Wiles,TaylorWiles}, and in the theorem of Breuil, Conrad, Diamond and Taylor for all elliptic curves over the rationals \cite{BCDT}. Serre's formulation of residual modularity \cite{Serre} and Ribet's level-lowering theorem \cite{Ribet} provide the framework for passing from a Frey curve to a modular form of prescribed level.

Bennett, Vatsal and Yazdani \cite{BVY} develop the level-lowering method used here. Their Theorem 1.3 excludes the coefficient prime 2. Lemmas 3.1 and 3.4 give absolute irreducibility and the residual representation needed for the present equation. We continue from that representation to an explicit determination of the relevant points over its degree-18 line-stabilizer field.

Freitas, Naskr\k{e}cki and Stoll \cite{FNS} combine modular methods and elliptic Chabauty for the generalized Fermat equation with exponents two, three and a prime. Their results for exponents eleven and thirteen provide related methodological context: the complete exponent-eleven result is conditional on GRH, while the exponent-thirteen results are partial. For the ternary equation considered here, we establish the class-group and descent inputs unconditionally and carry the local analysis through all relevant residue disks.

The local analysis belongs to the broader use of p-adic methods for determining rational points. Coleman developed an effective form of the classical Chabauty method \cite{Coleman}, while Bruin developed Chabauty methods using elliptic curves over number fields \cite{Bruin}. In the present calculation, rationality of the modular invariant supplies the equations on the local Mordell--Weil closure.

\paragraph{Modular reduction} A hypothetical solution defines a Frey curve. Level lowering and local comparison select a single residual representation at level 54. Its action on cyclic subgroups gives a line-stabilizer field of degree 18 and a point on the modular elliptic curve over that field. The Frey family forces the modular invariant to be rational and congruent to two modulo five.

\paragraph{Global arithmetic and the rank} Exact field and modular-map identities identify the objects on which the arithmetic is performed. An unconditional prime-ideal generation bound and exhaustive ideal relations prove that the class group is cyclic of order 16, giving a global squareclass space of dimension 14. Complete local Kummer images determine the Selmer groups. The lower pairing has rank six. On the residual two-dimensional quotient by the known point classes, the higher pairing has rank two and gives a Mordell--Weil rank bound of three. Three independent points attain this bound.

\paragraph{Local closure and the residue disks} Saturation at two and five shows that the subgroup generated by the explicit points and rational torsion has the same closure at five as the full Mordell--Weil group. This gives three formal parameters on each of the 18 relevant residue disks. Exact arithmetic modulo 25 excludes 15 disks. On each of the remaining three, an invertible derivative minor and integral power-series estimates prove uniqueness of the parameter whose modular invariant lies in the rational 5-adic field.

\paragraph{The final classification} The modular-point computation leaves three modular invariants. For the ternary Frey family, each gives a rational equation of degree four. Factoring these equations and imposing the seventeenth-power condition determines all primitive solutions.

Sections 2 and 3 carry out the modular reduction and identify the field and modular map. Section 4 establishes the class group and global squareclass basis. Sections 5 and 6 determine the Selmer groups and the rank bound. Section 7 proves the saturation and local-closure statements, then determines the relevant modular invariants. Section 8 completes the primitive classification and derives the stated consequences.

Throughout, squareclasses are written multiplicatively, with their linear
algebra performed over $\F_2$. All ranks of binary matrices mean ranks over
$\F_2$. For a number field $K$ and a finite set $S$ of finite primes, put
\[
K(S,2)=\{a\in K^*/K^{*2}:v_{\mathfrak q}(a)\equiv0\pmod2
\text{ for every }\mathfrak q\notin S\}.
\]
An elliptic-curve rank always means its algebraic Mordell--Weil rank.

\file{Supplementary_Data_Guide.txt} gives a section-by-section index to the exact algebraic data, local calculations and verification programs used in the proof. In particular, it locates the complete ideal enumeration, the Selmer and pairing data, and the residue-disk calculations.

\section{From a primitive solution to a modular point}\label{sec:modular}

\subsection{Primitivity and the prime three}

Let $X,Y,Z$ be nonzero integers with $\gcd(Y,Z)=1$ and $X^3=Y^{17}+2Z^{17}$. If $Y$ were even, $Z$ would be odd and the right-hand side would have $2$-adic valuation $1$, which is impossible for a cube. Hence $X$ and $Y$ are odd. The equation also gives
\[
\gcd(X,Y)=\gcd(X,2Z)=1.
\]
If $|YZ|=1$, direct substitution of the four pairs $Y,Z\in\{1,-1\}$ gives exactly the two solutions in Theorem~\ref{thm:ternary}. We may therefore suppose that $|YZ|>1$.

If $3\mid Z$, use the Bennett--Vatsal--Yazdani normalization $A=1$, $B=2$, $C=1$, $a=Y$, $b=Z$, $c=X$. If $3\mid Y$, interchange the two seventeenth-power terms, taking $A=2$, $B=1$, $a=Z$, $b=Y$ instead. In either case $3\nmid Aa$, while $v_3(Bb^{17})\ge17$. Corollary 3.3 and Lemma 3.4 of \cite{BVY} give residual level $6$. There is no weight-two cuspidal newform at this level, a contradiction. Thus
\[
3\nmid YZ.
\]

\subsection{The Frey curve and residual representation}

We use the Weierstrass invariant and reduction conventions of Silverman \cite{Silverman} and Tate's algorithm for the local reduction types \cite{Tate}.

Consider the Frey curve
\[
E_{X,Z}:\quad v^2+3Xuv+2Z^{17}v=u^3.
\]
Its invariants are
\[
\Delta(E_{X,Z})=216Y^{17}Z^{51},\qquad
c_4(E_{X,Z})=9X(9X^3-16Z^{17}).
\]
Since $X$ is odd, $c_4$ is a $2$-adic unit. The model is minimal at $2$, and
\[
v_2(\Delta)=3+51v_2(Z)\equiv3\pmod{17}.
\]
Its reduction modulo $2$ has nodal tangent cone $v(v+u)$, so the multiplicative reduction is split and the residual inertia action is nontrivial.

Choose $\varepsilon\in\{1,-1\}$ such that $2\varepsilon Z^{17}\equiv1\pmod3$. The normalization
\[
(A,B,C;a,b,c)=(\varepsilon,2\varepsilon,\varepsilon;Y,Z,X)
\]
satisfies the hypotheses of \cite[Lemmas 3.1 and 3.4]{BVY}. Its Frey model has coefficients $a_1=3\varepsilon X$ and $a_3=2\varepsilon Z^{17}$, and is isomorphic over $\mathbb Q$ to $E_{X,Z}$ by changing the sign of $v$ when necessary. Because $17\nmid ABC$, the residual representation arises at level $18$ or $54$. Level $18$ has no weight-two cuspidal newform. The two rational newforms at level $54$ are represented by the curves $E_a$ and $E_b$ listed below.

The two rational newforms at level $54$ are represented by
\[
E_a=[1,-1,0,12,8],\qquad E_b=[1,-1,1,1,-1],
\]
in the usual five-coefficient Weierstrass notation. Their coefficients at
$2$ and $5$ are listed in Table~\ref{tab:1}.
\begin{table}[htbp]
\centering
\caption{Fourier coefficients of the two level-54 newforms.}\label{tab:1}
\begin{tabular}{lrr}
\toprule
 & $a_2$ & $a_5$\\
\midrule
$E_a$ & $-1$ & $3$\\
$E_b$ & $1$ & $-3$\\
\bottomrule
\end{tabular}
\end{table}
Let
\[
E_*:\quad v^2+3uv+2v=u^3.
\]
It is isomorphic to $E_b$, and
\[
(c_4,c_6,\Delta,j)(E_*)=(-63,351,-216,9261/8).
\]


On the inertia-invariant line at $2$, the Frobenius eigenvalue for $E_{X,Z}[17]$ is $2$, while it is $-2$ for the nonsplit multiplicative curve $E_a$. These differ modulo $17$, and hence $E_a$ is excluded. As $E_b$ is isomorphic to the reference curve $E_*$, we obtain
\[
E_{X,Z}[17]\simeq E_*[17]
\]
as representations of $G_{\mathbb Q}$ over $\mathbb F_{17}$. 
\subsection{The line-stabilizer field}

At $2$ the reference representation contains a nontrivial transvection.
At $5$ its characteristic polynomial is $T^2+3T+5$, with discriminant
$6$ modulo $17$, a nonsquare. Hence it is irreducible. Some group element
moves the fixed line of the transvection. In a basis consisting of the two
fixed lines, the resulting two transvections are nontrivial upper and
lower unipotent matrices. Since $\F_{17}$ is the prime field, their
powers include every upper and lower elementary unipotent matrix. They
generate $\operatorname{SL}_2(\F_{17})$. The determinant is the
surjective cyclotomic character, so the full image is
$\operatorname{GL}_2(\F_{17})$.

Fix a line $C_*\subset E_*[17]$, and let $L$ be its field of definition.
The action on the $18$ lines is transitive, so $[L:\Q]=18$.
Complex conjugation has two fixed lines and eight two-cycles, giving
signature $(2,8)$. 

Transporting the fixed line in $E_*[17]$ gives a point $Q\in X_0(17)(L)$ with $j(Q)=j(E_{X,Z})\in\mathbb Q$.

\subsection{The local condition at five}

If $5\mid YZ$, then $5\nmid X$ and $c_4$ is a $5$-adic unit. The Frey curve has multiplicative reduction, with minimal discriminant valuation divisible by $17$. Its unramified residual trace at $5$ is therefore $\pm6$, whereas the reference representation has trace $-3$ modulo $17$. This is impossible. Thus $5\nmid YZ$, the Frey curve has good reduction at $5$, and the Hasse bound together with the residual congruence gives $a_5(E_{X,Z})=-3$.

For completeness, the finite calculation of the required congruence uses the reductions $r=X\bmod5$ and $z=Z\bmod5$. For each $r\in\mathbb F_5$ and $z\in\mathbb F_5^*$, count the affine solutions of
\[
v^2+3ruv+2zv=u^3
\]
and add the point at infinity, discarding the singular pairs. Here $Z^{17}\equiv z\pmod5$. The pairs of trace $-3$ are listed in the accompanying verification data; each has $j=2$ in $\mathbb F_5$. Consequently
\[
j(E_{X,Z})\in\mathbb Z_5,\qquad j(E_{X,Z})\equiv2\pmod5.
\]
Proposition~\ref{prop:rationalj} now applies and gives
\[
j(E_{X,Z})\in\left\{\frac{9261}{8},-\frac{132651}{2},-\frac{1167051}{512}\right\}.
\]


\section{Explicit field and modular-curve models}\label{sec:field}

\subsection{A small line invariant and the two field models}

The change of variables $x=4u+3$, $y=8v+12u+8$ takes $E_*$ to
$E':y^2=x^3+21x-26$. If $C_*=\langle P\rangle$, define
\[
\tau=\frac14\sum_{k=1}^{8}x(kP).
\]
It is fixed by the line stabilizer. In
$\Q[x]/(\psi_{17}(x))$, where $\psi_{17}$ has degree $144$,
division-polynomial arithmetic gives the following polynomial for $\tau$:
\begin{align*}
F(T)={}&T^{18}+4284T^{16}-169728T^{15}-7197120T^{14}
 +99800064T^{13}\\
&-174311472T^{12}-15568809984T^{11}-79963790256T^{10}
 +700279075600T^9\\
&-3717675394776T^8+2141406520344T^7+4041758761800T^6
 +63456301182456T^5\\
&-192629564298360T^4-170217634882776T^3
 -45116978926824T^2\\
&-19005735615048T-39301585356176.
\end{align*}
Its reduction modulo $5$ has increasing coefficient list
\[
(4,2,1,4,0,1,0,4,4,0,4,1,3,4,0,2,4,0,1).
\]
The Rabin irreducibility certificate checks $T^{5^{18}}=T$ modulo $F$
and $\gcd(T^{5^d}-T,F)=1$ for $d=6,9$. Thus $F$ is irreducible.
Since $\Q(\tau)\subseteq L$ and both degrees are $18$, it defines the
required line field.

Multiplication-map arithmetic gives $F(\tau)=0$ in
the $144$-dimensional division algebra. Since $F$ is irreducible, the
characteristic polynomial of multiplication by $\tau$ is $F(T)^8$.
V\'elu's formulas compute the quotient curve
$E'/C_*$ and its $j$-invariant as explicit elements of $\Q[T]/(F)$.
Substitution in the classical modular polynomial $\Phi_{17}$ gives zero.
The modular polynomial has $172$ symmetric coefficient records; its identification is proved below.

For the global ideal and descent computations we use $L=\Q(\alpha)$ with
\begin{align*}
g(Y)={}&Y^{18}-6Y^{17}+136Y^{15}-612Y^{14}-3060Y^{13}
 +46444Y^{12}\\
&-208998Y^{11}+417996Y^{10}-376278Y^9+1701258Y^8
 -12150648Y^7\\
&+41782464Y^6-84353490Y^5+109537290Y^4-93550626Y^3\\
&+51213384Y^2-16413030Y+2359296.
\end{align*}
The explicit element $\beta\in\Q[Y]/(g)$ satisfies $F(\beta)=0$ and
has minimal polynomial $F$, giving the field isomorphism. The field has signature and discriminant
\begin{equation}\label{eq:disc}
(r_1,r_2)=(2,8),\qquad
D_L=2^{16}3^{20}17^{17}
=189032591590977297735036384246300672.
\end{equation}

\subsection{The genus-one model and its isogenies}

The modular interpretation of cyclic isogenies and modular curves follows Diamond and Shurman \cite{DS}. The quotient-curve calculations use V\'elu's explicit isogeny formulas \cite{Velu}.

Write
\begin{equation}\label{eq:A}
A=X_0(17):\quad y^2+xy+y=x^3-x^2-x-14.
\end{equation}
The $j$-map used throughout is
\begin{equation}\label{eq:jmap}
j(x,y)=\frac{a(x)+b(x)y}{(x-7)^{17}},
\end{equation}
where the integer coefficients of $a$ and $b$ are displayed in Appendix~\ref{app:jmap}. Their degrees are $17$ and $16$.
The map has a pole of order $1$ at the origin $O$ and a pole of order
$17$ at $T=(7,13)$. The expression at $-T$ is removable.

Two useful isogenous models are
\[
A_0:\ V^2=U(U^2+30U+289),\qquad
F_0:\ w^2=t(t-16)(t+1).
\]
Here $U=4x-11$, $V=8y+4x+4$ gives an isomorphism $A\simeq A_0$,
and the degree-$2$ map from $A_0$ to $F_0$ is
\[
t=\frac{V^2}{4U^2},\qquad
w=\frac{V(289-U^2)}{8U^2}.
\]
The formulas extend across their apparent exceptional points. A further
degree-$2$ isogeny $\phi:F_0\to C$ has kernel $\langle(-1,0)\rangle$,
where
\[
C:\ v^2=u(u^2+36u+256).
\]
All these elliptic curves have the same rank over $L$.


% BEGIN field_bridge_section.tex
\subsection{The field and modular map}
\label{subsec:field-bridge-independent}

We now give exact identities identifying the field and the modular $j$-map.

\paragraph{The kernel field and the change of generator.}
Division-polynomial arithmetic on $E':y^2=x^3+21x-26$ gives the squarefree degree-$144$ polynomial $\psi_{17}$. For $1\leq k\leq8$, the rational multiplication maps $x([k]P)$ have denominators coprime to $\psi_{17}$. In the resulting \'etale algebra, their sum $S=\sum_{k=1}^{8}x([k]P)$ satisfies $F(S/4)=0$.
Since $F$ is irreducible of degree $18$, the $18$ cyclic subgroups have
$18$ distinct invariant values. Indeed, the sum is constant on each
subgroup, there are only $18$ such subgroups, and the irreducible
annihilating polynomial forces all $18$ conjugates to occur. Thus the
stabilizer of $S/4$ is precisely a line stabilizer.

The reduced polynomial $g$ is related to $F$ by an explicit element $\beta\in\Q[Y]/(g)$ satisfying
\[
F(\beta)=0,\qquad \operatorname{minpoly}_{\Q}(\beta)=F.
\]
These identities establish the field isomorphism used in the ideal and descent calculations. The integral basis and maximal-order data are included in the supplementary material.

\paragraph{An exact certificate for $\Phi_{17}$.}
We use the coefficient polynomial for $\Phi_{17}$ in \cite{Sutherland}. The prime-level coefficient data are computed by the isogeny-volcano algorithm of Br\"oker, Lauter and Sutherland \cite{BLS}. The following integer power-series calculation identifies it directly. Put
\[
E_4(q)=1+240\sum_{n\geq1}\sigma_3(n)q^n,\qquad
E_6(q)=1-504\sum_{n\geq1}\sigma_5(n)q^n,
\]
\[
\Delta(q)=\frac{E_4(q)^3-E_6(q)^2}{1728},\qquad
j(q)=\frac{E_4(q)^3}{\Delta(q)}.
\]
For the candidate polynomial $P(U,V)$, every coefficient of
$P(j(q),j(q^{17}))$ from $q^{-306}$ through $q^{306}$ vanishes exactly.
For every monomial $U^iV^j$ present, the computation also checks
\[
i+17j\leq306,\qquad 17i+j\leq306.
\]
The function $P(j(z),j(17z))$ on the standard modular curve $X_0(17)$
has no poles away from its two cusps. At those cusps, the pole orders
of the two arguments are $(1,17)$ and $(17,1)$, respectively. Hence its
pole order at either cusp is at most $306$. The certified expansion
gives a zero of order at least $307$ at infinity. If the function were
nonzero, its sole remaining pole would have order at most $306$,
contradicting the degree-zero property of a principal divisor. The
function is therefore identically zero. Monicity, degree $18$, and
$[\operatorname{SL}_2(\Z):\Gamma_0(17)]=18$ identify $P$ with the
classical modular polynomial. The integer power-series calculation therefore certifies the coefficient polynomial exactly.

\paragraph{Identification of the modular $j$-map.}
Temporarily regard the curve in \eqref{eq:A} simply as an explicitly
given genus-one curve, and set $T=(7,13)$ and $w(P)=T-P$.
The group law verifies this involution as a rational self-map. Exact
substitution into \eqref{eq:jmap} gives
\[
j(w(x,y))=\frac{c(x)+d(x)y}{x-7},
\]
where
\begin{align*}
c(x)={}&-8x^9+28x^8+269x^7-350x^6-652x^5-3704x^4\\
&+1776x^3+4960x^2+8640x-2688,\\
d(x)={}&x^8+21x^7-163x^6-316x^5+372x^4+1264x^3\\
&+2288x^2-1344x-832.
\end{align*}
In its exact function field, reduction modulo
$y^2+(x+1)y-(x^3-x^2-x-14)$ verifies
\begin{equation}\label{eq:fricke-modular-identity}
\Phi_{17}\bigl(j(P),j(T-P)\bigr)=0.
\end{equation}
After clearing $(x-7)^{306}$, the polynomial coefficients of both $1$ and $y$ vanish identically.

The pole order of $j$ at $O$ is $1$, since the numerator has pole order
$35$ and the denominator has pole order $34$. At $T$, the numerator is
nonzero, giving pole order $17$. At $-T=(7,-21)$, the exact polynomial
\[
a(x)^2-(x+1)a(x)b(x)-(x^3-x^2-x-14)b(x)^2
\]
has valuation $17$ at $x-7$. The conjugate numerator is nonzero there,
so the apparent denominator cancels completely. There are no other
poles, and therefore $\deg j=18$.
The classical polynomial $\Phi_{17}(j,V)$ is irreducible of degree
$18$ over $\Q(j)$. Equation~\eqref{eq:fricke-modular-identity} and the
degree calculation imply
\[
\Q(A)=\Q\bigl(j(P),j(T-P)\bigr).
\]
This identifies the displayed curve and map with the canonical
$X_0(17)$ function field. 

\paragraph{The exact V\'elu-to-fibre bridge.}
From the multiplication maps we compute
$S_2=\sum_{k=1}^{8}x([k]P)^2$ and
$S_3=\sum_{k=1}^{8}x([k]P)^3$. Their expressions in $\Q[S/4]$ are given in the supplementary material. V\'elu's quotient has coefficients
\[
A'=-79\cdot21-30S_2,\qquad
B'=-223(-26)-42\cdot21S-70S_3,
\]
and its $j$-invariant equals
$6912A'^3/(4A'^3+27B'^2)$.
Eliminating $y$ from $j(x,y)=9261/8$ and \eqref{eq:A} gives an irreducible degree-$18$ fibre polynomial,
including the exact cancellation of the spurious factor
$(x-7)^{17}$. The specified point $P_1\in A(L)$ satisfies the curve
equation and $j(P_1)=9261/8$, and its $x$-coordinate generates $L$.
Finally, in the same reduced field, the calculation checks
\[
j_{\mathrm{quot}}(\beta)=j(T-P_1)
\]
exactly; this common element also has degree $18$.
Thus the subgroup invariant, the reduced field, the V\'elu quotient,
and the modular fibre are connected by exact identities.

% END field_bridge_section.tex


\section{The unconditional global squareclass space}\label{sec:global}

\subsection{An ideal class of exact order $16$}

The integral-basis, ideal and unit computations use standard computational algebraic number theory as developed by Cohen \cite{Cohen1993,Cohen2000}. The following witnesses determine the ideal classes and squareclasses needed for the descent.

The exact certificate specifies a fractional ideal $\mathfrak P$, an element
$\gamma$, and ten units $U=(-1,\epsilon_1,\ldots,\epsilon_9)$ such that
\[
\mathfrak P^{16}=(\gamma).
\]
For every displayed unit, its principal ideal is the unit ideal and its norm
has absolute value $1$. Eleven finite quadratic residue characters give
an invertible binary matrix on
$[\gamma,-1,\epsilon_1,\ldots,\epsilon_9]$.

Dirichlet's theorem gives $\dim_{\F_2}\OO_L^*/\OO_L^{*2}=10$.
Consequently the ten displayed units are a complete basis modulo squares.
If the order of $[\mathfrak P]$ divided $8$, an identity
$\mathfrak P^8=(\beta)$ would imply $\gamma=u\beta^2$ for a unit $u$.
The invertible eleven-character matrix excludes this. Hence
\begin{equation}\label{eq:Porder}
\operatorname{ord}([\mathfrak P])=16.
\end{equation}
The unit squareclasses and the order of the ideal class are therefore determined by exact identities and residue characters.

For each of the $112$ factor-base prime ideals $\mathfrak q_i$,
the certificate provides an element $\gamma_i$ and an exponent $e_i$ with
\[
\mathfrak q_i=(\gamma_i)\mathfrak P^{e_i}.
\]
Canonical ideal HNF equality checks every identity. Thus every factor-base class lies in
$H=\langle[\mathfrak P]\rangle\simeq\Z/16\Z$.

\subsection{Exhaustive class generation}

We use the unconditional prime-ideal criterion of
Belabas--Diaz y Diaz--Friedman, Corollary~5.2 of \cite{BDDF}.
For a degree-$n$ field of discriminant $D$ and $r_1$ real places, put
\[
\mathcal B(T)=\log D-n\left(\gamma+\log(4\pi)
 -\frac{\pi^2}{4\log T}\right)
 -r_1\left(1-\frac{\log4}{\log T}\right),
\]
where $\gamma$ is Euler's constant. The criterion states that if
\begin{equation}\label{eq:analytic-generation}
4\sum_{(\mathfrak q,m):\,N\mathfrak q^m<T}
 \frac{\log N\mathfrak q}{1+(N\mathfrak q)^m}
 \left(1-\frac{m\log N\mathfrak q}{\log T}\right)
 >\mathcal B(T),
\end{equation}
then the prime ideals of norm strictly less than $T$ generate the ideal
class group. This is the unconditional result in Section~5 of that paper;
it does not assume the Riemann hypothesis for any $L$-function.

Here we take
\begin{equation}\label{eq:B}
T=10\,000\,001,\qquad B=T-1=10\,000\,000.
\end{equation}
Every term on the left of \eqref{eq:analytic-generation} is positive.
We may therefore replace the sum by the partial sum over the $5068$
distinct prime ideals of norm at most $50000$, retaining for each ideal
all positive integers $m$ with $(N\mathfrak q)^m<T$.
The exact census has $2659$ rows $(p,f,\text{multiplicity})$ and gives
$5630$ terms after multiplicities are included. Ramification indices do
not multiply the terms. Using the prime-ideal decompositions in the maximal order, $256$-bit outward-rounded real ball arithmetic gives
\[
26.23442619<\text{partial left side}<26.23442620,
\]
\[
26.20641426<\mathcal B(T)<26.20641427.
\]
In particular the strict inequality has margin greater than $0.02801192$.
This proves the generation bound directly from exact prime-ideal data and
certified elementary real functions.

The positive partial sum establishes the analytic inequality, whose conclusion is generation by prime ideals below the larger bound. We next express every prime ideal up to that bound in the subgroup generated by the specified ideal class.

The prime ideals of norm at most $B$ are enumerated over $20$ adjacent rational-prime intervals, each of width $500000$ except the last. The norm bound is $B$ throughout. Complete prime decomposition in the first interval includes every residue degree with prime-ideal norm at most $B$. The enumeration contains $664579$ rational primes and $664562$ prime ideals; the first interval contains $41681$ prime ideals, including $156$ of residue degree greater than one.

In the remaining intervals, $p>\sqrt B$, so every relevant prime ideal has residue degree $1$. These primes divide neither the power-basis index nor the discriminant. The roots of $g$ modulo $p$ therefore enumerate the relevant prime ideals. For a root $r$ and the integral basis $1=b_1,b_2,\ldots,b_{18}$, the corresponding ideal is the kernel of
\[
\OO_L\longrightarrow\F_p,\qquad b_i\longmapsto b_i(r).
\]
The columns $pe_1,e_2-b_2(r)e_1,\ldots,e_{18}-b_{18}(r)e_1$ give an integral ideal basis of determinant $p$; reducing the first-row entries to canonical residues gives its HNF.

The factor base is enlarged to the $5068$ prime ideals of norm at most $50000$. Principal-ideal relations from the first interval express their classes in $H$. For each target ideal $I$, lattice reduction supplies elements $m\in I$. The exact quotient $|N(m)|/N(I)$ and the valuations at the factor-base primes determine the relation. Equality of norms after removing these valuations establishes the full prime-ideal support, and hence expresses $[I]$ in $H$.

Norm evaluation uses outward-rounded embedding intervals with denominator $2^{48}$, covering two real embeddings and eight conjugate complex pairs. For the bounded coefficients used here, interval linear forms are evaluated with signed $128$-bit integers and a bound below $2^{86}$; products use arbitrary precision integers. A quotient norm is determined from an interval only when that interval contains a unique possible integer; otherwise the algebraic norm is evaluated exactly.

The resulting principal-ideal relations place every prime ideal of norm at most $B$ in $H$. The generation criterion \eqref{eq:analytic-generation}, together with the exact order \eqref{eq:Porder}, therefore gives
\begin{equation}\label{eq:class}
\Cl(\OO_L)=H\simeq\Z/16\Z.
\end{equation}

\subsection{Fourteen complete global squareclasses}

Let $S=\{\mathfrak p_1,\mathfrak p_2,\mathfrak p_{17}\}$ be the two
primes over $2$ and the unique prime over $17$. Their $(e,f)$ values are
$(17,1),(1,1),(18,1)$, respectively. Exact principal-ideal witnesses give
\[
[\mathfrak p_1]=2[\mathfrak P],\qquad
[\mathfrak p_2]=14[\mathfrak P],\qquad
[\mathfrak p_{17}]=0.
\]
It follows from \eqref{eq:class} that the $S$-class group is $\Z/2\Z$.
The exact sequence
\[
0\longrightarrow\OO_{L,S}^*/\OO_{L,S}^{*2}
\longrightarrow L(S,2)\longrightarrow\Cl(\OO_{L,S})[2]
\longrightarrow0
\]
therefore gives
\begin{equation}\label{eq:Gdim}
\dim L(S,2)=(r_1+r_2+|S|)+1=13+1=14.
\end{equation}

Fourteen exact elements $G_1,\ldots,G_{14}$ are supplied. Their principal
ideals are
\[
(G_i)=\OO_L\ (1\le i\le10),\quad
(G_{11})=\mathfrak p_1^8,\quad (G_{12})=\mathfrak p_1\mathfrak p_2,
\]
\[
(G_{13})=\mathfrak p_{17},\qquad
(G_{14})=\mathfrak p_1^3\mathfrak q_{719}^{2}.
\]
Thus they all belong to $L(S,2)$. Fourteen independently constructed
quadratic residue characters, over the rational primes
\[
11,11,13,13,19,19,37,37,41,41,43,43,47,67,
\]
give a $14\times14$ matrix with determinant $1$ in $\F_2$.
Together with \eqref{eq:Gdim} this proves that the $G_i$ form a complete
basis. The resulting squareclass basis is obtained directly from these ideal and character calculations.

\section{Local descent and the lower pairing}\label{sec:lower}

\subsection{The complete local conditions}

The Selmer-group construction uses the cohomological framework of Cassels \cite{CasselsIII} and the explicit descent methods described by Schaefer and Stoll \cite{SchaeferStoll}. Algorithms for rank computations over number fields are developed by Simon \cite{Simon}. The global squareclasses and local Kummer conditions for the present curves are given below.

For $F_0:w^2=t(t-16)(t+1)$, the full $2$-Kummer map has coordinates
\[
(t,t-16,t+1)\in(L^*/L^{*2})^3,
\]
with the usual limiting values at $2$-torsion points. Their product is a
square, so two coordinates specify the class. Outside $S$ the curve has
good reduction and odd residue characteristic, and its local Kummer image
is unramified. Consequently the global search space is exactly
$L(S,2)^2$, of dimension $28$.

At the degree-$17$ and degree-$1$ dyadic completions, the local
squareclass dimensions are $19$ and $3$, and the full Kummer-image
dimensions are also $19$ and $3$. At the prime over $17$ the Kummer
image has dimension $2$. These are complete images: explicit local point
classes attain the general local dimension bounds. At the real places,
the allowed signs of the first two coordinates are $(+,+)$ and $(-,-)$.
Their exact values are certified by rational root isolation and interval
Horner arithmetic.

Expressing the local images in the basis $G_i$ gives the constraint ranks in Table~\ref{tab:3}:
\begin{table}[htbp]
\centering\caption{Ranks of the local Kummer constraints.}\label{tab:3}
\begin{tabular}{lr}
\toprule
Conditions imposed&Rank in $28$ variables\\\midrule
First dyadic place&13\\
Both dyadic places&14\\
Dyadic places and the prime over $17$&15\\
All finite and real conditions&15\\\bottomrule
\end{tabular}
\end{table}
The complete constraint matrix has size $26\times28$; its binary entries are included in the supplementary data. Hence
\[
\dim\Sel^{(2)}(F_0/L)=28-15=13.
\]

The diagonal classes $(a,a,1)$ form the initial isogeny Selmer group
$S_1$ for $\phi:F_0\to C$. Substituting equal coordinate vectors in
the constraint matrix gives a space of dimension $11$. The exact
representatives $A_1,\ldots,A_{11}$ in the supplementary data form its basis:
for each $A_i$, its expression in the $G_j$ is accompanied by an exact
square root of the quotient, and the eleven coordinate columns have
rank $11$.

The opposite initial group $S'_1$ is computed separately. The local
Kummer map for the dual isogeny on $F_0$ is the squareclass of $t+1$;
it is the sum of the two binary coordinates above. Projecting each
\emph{complete local image}, then imposing its conditions on all
fourteen global squareclasses, gives dimension $2$. This computes the opposite Selmer group from the complete local images.

Both basis classes of $S'_1$ are represented by actual points. They are
$17$ and $u_1=4x(P_1)-11$, where $P_1$ is the point defined in
Section~\ref{sec:points}. The first comes from $(16,0)\in F_0(L)$.
For the image of $P_1$ under $A_0\to F_0$, the identity
\[
t+1=\frac{(U+17)^2}{4U}
\]
gives the second. Exact squareclass calculations prove their independence. Thus $S'_1$ is generated by Mordell--Weil classes.

\subsection{The lower Cassels matrix}

The pairing is considered in Cassels' duality theory for curves of genus one \cite{CasselsIV}. Its local factors are evaluated on the explicit covering curves below.

For a diagonal class $a$, use the covering
\[
a z_1^2=z_3^2-z_0^2,\qquad
a z_2^2=z_3^2-17z_0^2.
\]
The certificate gives a nonzero exact conic point $(r,s,h)$ satisfying
$r^2-a s^2=17h^2$. Take the tangent forms
\[
L_1=as z_2-rz_3+17h z_0,\qquad L_2=z_3-z_0.
\]
The lower multiplicative pairing with a second diagonal class $b$ is
\begin{equation}\label{eq:lowerpairing}
\prod_v(L_1(P_v)L_2(P_v),b)_v,
\end{equation}
for suitable local covering points $P_v$. This is the Cassels--Tate
pairing by the formula and comparison theorem in the preprint of Shukla and Stoll (2023, arXiv:2302.01640).

Every $A_i$ has even valuation at every finite place and is a square
at the prime above $17$. If $(a)=\mathfrak A^2$, the good odd-place
contribution can be computed from the fractional content ideal
\[
\mathfrak C=(r,h)+\mathfrak A(s).
\]
Indeed, after the square rescaling which makes $a$ a unit, divide the
conic point by a uniformizer to the power $v(\mathfrak C)$.
A primitive tangent cuts a smooth conic in twice a point; its evaluation
has even valuation. The product in \eqref{eq:lowerpairing} therefore has
valuation congruent to $v(\mathfrak C)$ modulo $2$. Since $b$ has even
valuation, the factor is
$(\pi,b)_v^{v(\mathfrak C)}$. Complete ideal factorization finds every
possible odd contribution. Real signs are calculated exactly, and the
factor above $17$ is $1$.

At the large dyadic place, independent exact norms $x^2-b y^2$ span an
$18$-dimensional hyperplane in the $19$-dimensional squareclass space.
Its unique nonzero annihilator is the Hilbert character for $b$.
Applying this construction to all eleven $b$ values determines the dyadic characters. The local points and tangent
squareclasses are certified by Hensel inequalities and valuation error
bounds. Combining the dyadic, good odd and real factors gives the matrix below.

Writing the pairing as $(-1)^{M_{ij}}$, the result is
\begin{equation}\label{eq:lowerM}
M=\begin{pmatrix}
0&0&0&0&0&0&0&0&0&0&0\\
0&0&0&0&0&1&0&1&1&0&1\\
0&0&0&0&1&0&1&0&1&0&1\\
0&0&0&0&1&0&1&1&1&0&1\\
0&0&1&1&0&1&0&1&1&0&1\\
0&1&0&0&1&0&1&0&1&0&0\\
0&0&1&1&0&1&0&1&1&0&1\\
0&1&0&1&1&0&1&0&0&1&1\\
0&1&1&1&1&1&1&0&0&0&1\\
0&0&0&0&0&0&0&1&0&0&0\\
0&1&1&1&1&0&1&1&1&0&0
\end{pmatrix},\qquad\rank M=6.
\end{equation}
It is alternating. Its radical has dimension $5$, with representatives
\[
R_1=-1,\quad R_2=A_5A_7,\quad
R_3=A_3A_4A_5A_6A_9,\quad R_4=A_3A_4A_{10},
\]
\[
R_5=A_2A_3A_4A_5A_6A_{11}.
\]
The known point subspace has radical-coordinate basis
\[
(1,0,0,0,0),\quad (0,1,0,1,1),\quad (1,0,1,0,1).
\]
Exact square identities identify these three radical directions with point classes. The higher pairing in Section 6 eliminates the remaining two directions and gives the rank bound.


% BEGIN higher_pairing_section.tex
\section{The higher isogeny pairing and the rank bound}
\label{sec:higher-pairing}

The global squareclass basis and complete local images from Section 5 provide the input for the higher descent. We describe the Selmer filtration and then compute a nonzero higher pairing on the two-dimensional quotient of the lower-pairing radical by the known point classes.

In this section abbreviate the curve $F_0$ by $F$ (the polynomial
$F(T)$ from the field construction is a different object). Put
\[
 F:\ w^2=t(t-16)(t+1),\qquad
 C:\ v^2=u(u^2+36u+256).
\]
There is a degree-two isogeny $\phi:F\longrightarrow C$ whose kernel is
generated by $(-1,0)$; write $\psi:C\longrightarrow F$ for its dual.
Indeed, setting $r=t+1$ writes $F$ as
$w^2=r(r^2-18r+17)$, and the usual two-isogeny quotient has equation $C$.
The ranks of these curves over $L$ agree with the rank of the modular
elliptic curve used in the preceding sections.

\subsection{The Selmer filtration and rank bound}

We use the notation of Fisher~\cite{Fisher}, Section~2 and Theorem~2.1.
The first four terms of the filtration on the $\phi$ side are
\[
\begin{aligned}
 S_1&=\operatorname{Sel}^{\phi}(F/L),\\
 S_2&=\operatorname{im}\bigl(\operatorname{Sel}^{2}(C/L)
                    \longrightarrow S_1\bigr),\\
 S_3&=\operatorname{im}\bigl(\operatorname{Sel}^{2\phi}(F/L)
                    \longrightarrow S_1\bigr),\\
 S_4&=\operatorname{im}\bigl(\operatorname{Sel}^{4}(C/L)
                    \longrightarrow S_1\bigr).
\end{aligned}
\]
The primed filtration is obtained by exchanging $F,C$ and $\phi,\psi$.
Here $2\phi=[2]_C\circ\phi$ has degree eight, while $[4]_C$ has degree
sixteen.  The index counts successive two-isogenies; it does not denote
multiplication by $2^m$.

For odd $m$, Fisher's pairing $\Theta_m$ is alternating on $S_m$, with
kernel $S_{m+1}$.  For even $m$ it is a pairing
$S_m\times S'_m\longrightarrow\mathbb F_2$, with left and right kernels
$S_{m+1},S'_{m+1}$.  We also use its multiplicative convention, in which
the nonzero value is $-1$.

Let
\[
 Q=C(L)/\phi F(L),\qquad Q'=F(L)/\psi C(L).
\]
For every $m$ there are inclusions $Q\subseteq S_m$ and $Q'\subseteq S'_m$.
Both rational isogeny kernels have dimension one, so the elementary
Mordell--Weil exact sequence gives
\begin{equation}\label{eq:rank-isogeny-quotients}
 \operatorname{rank}F(L)=\dim Q+\dim Q'-2
 \leq \dim S_m+\dim S'_n-2
 \quad(m,n\geq1).
\end{equation}
Applying these filtration inclusions to $S_4$ and $S'_4$ in \eqref{eq:rank-isogeny-quotients} gives the rank bound used below.

The complete $\phi$-Selmer space has dimension eleven.  Its basis classes
are represented in the full two-descent by $(A_i,A_i,1)$,
$1\leq i\leq11$.  The lower Cassels matrix has rank six, hence
\[
 \dim S_2=11-6=5.
\]
The initial opposite Selmer group computed in Section~\ref{sec:lower} has dimension two, with basis
\[
 S'_1=\langle 17,u_1\rangle,\qquad u_1=4x(P_1)-11.
\]
Both classes are represented by actual points of $F(L)$.  The class $17$
is the Kummer image of $(16,0)$.  For the other class, put
\[
 v_1=8y(P_1)+4x(P_1)+4.
\]
The exact point identity is
$v_1^2=u_1(u_1^2+30u_1+289)$.  Its image under the dual two-isogeny from
this auxiliary curve to $F$ is
\[
 t_1=\frac{u_1+30+289/u_1}{4},\qquad
 w_1=\frac{v_1(1-289/u_1^2)}8.
\]
It satisfies
\[
 w_1^2=t_1(t_1-16)(t_1+1),\qquad
 4u_1(t_1+1)=(u_1+17)^2.
\]
Consequently its opposite Kummer class is $u_1$. The local computation of $S'_1$ projects the complete local full-two-Kummer images by the third coordinate $t+1$. The projected image dimensions at the two dyadic places are $1,1$, the image above $17$ is the full two-dimensional squareclass space, and the real condition is positivity. Imposing these conditions on the global squareclass space gives the two classes exhibited above. This determines the initial group $S'_1$ before passing to its filtration term $S'_2$.

It follows that $Q'=S'_1$, and therefore $S'_m=Q'$ for every $m$.
The right-hand argument of $\Theta_2$ is always an actual
Mordell--Weil class, so $\Theta_2$ vanishes and
\[
 S_3=S_2,\qquad \dim S_3=5.
\]
In the radical basis
\[
\begin{aligned}
 R_1&=-1,& R_2&=A_5A_7,& R_3&=A_3A_4A_5A_6A_9,\\
 R_4&=A_3A_4A_{10},&
 R_5&=A_2A_3A_4A_5A_6A_{11},
\end{aligned}
\]
three independent actual Mordell--Weil classes have coordinates
\[
 (1,0,0,0,0),\quad(0,1,0,1,1),\quad(1,0,1,0,1).
\]
Write $W$ for their span.  Such classes lie in every filtration term
and in the radical of every higher pairing.
Their actual-point origin is also checked directly.  

The first class,
$-1$, is represented by $(-16,-32)\in C(L)$.  For each of the other
two specified points $R=(t,w)\in F(L)$, the certificate supplies a square
$s^2=t+1$.  The point
\[
 u_R=2(t-8+w/s),\qquad v_R=2u_Rs
\]
lies on $C$ and satisfies
$u_R=t(1+w/(st))^2$.
The square identity for $t$ therefore proves that its
$\phi$-Kummer class has precisely the displayed radical coordinates.
Exact point and square identities establish these representatives and their independence.

Choose
\[
 a=-17,\qquad b=-A_2A_3A_4A_6A_7A_{11},\qquad
 [b]=[-R_2R_5]\quad\text{in }L^*/L^{*2}.
\]
Their radical coordinates are respectively
\[
 (1,0,1,1,0),\qquad(1,1,0,0,1).
\]
The five displayed vectors have rank five over $\mathbb F_2$.
In particular, $a,b$ form a basis of $S_3/W$.  The square-equivalence
identifying $a$ with the specified product of the $A_i$ is checked by
an exact square root.  We shall prove
\begin{equation}\label{eq:theta3-nonzero}
 \Theta_3(a,b)=-1.
\end{equation}
Since the pairing is alternating and $W$ is in its radical,
its induced matrix on $S_3/W$ is
\[
 \begin{pmatrix}0&1\\1&0\end{pmatrix}.
\]
Thus its rank is exactly two, its kernel $S_4$ has dimension three,
and \eqref{eq:rank-isogeny-quotients} gives
\begin{equation}\label{eq:higher-rank-three}
 \operatorname{rank}F(L)\leq 3+2-2=3.
\end{equation}
The three independent points constructed in Section~\ref{sec:points} give equality.
Alternatively, the dimensions of the actual point quotient spaces $W$ and $Q'$ give the reverse inequality in
\eqref{eq:rank-isogeny-quotients}.

\subsection{The covering chain and its twists}

The diagonal covering attached to $a=-17$ is
\[
 \Gamma:\quad -17z_1^2=z_3^2-z_0^2,\qquad
                    -17z_2^2=z_3^2-17z_0^2.
\]
It has Jacobian $F$.  The map to the full two-descent target is
$t=z_3^2/z_0^2-1$, $w=-17z_1z_2z_3/z_0^3$.
The quotient by simultaneous sign change of $z_1,z_2$ is the
$\phi$-covering $\Gamma\longrightarrow C$.  For example its affine
coordinates can be written
\[
 u=\frac{2(z_3^2-9z_0^2-17z_1z_2)}{z_0^2},\qquad
 v=\frac{2uz_3}{z_0}.
\]
On $\Gamma$ the first coordinate also equals
$u=-17(z_1+z_2)^2/z_0^2$.  Thus the Kummer class of this
$\phi$-covering is exactly $a=-17$, with the stated sign.

Use the pushout and first twist
\[
 f_\Gamma=17(z_0-z_2)(z_3-z_0),\qquad \xi=-2.
\]
The first lift is
\[
 D:\quad
 D_1=17c^2+d^2-2\theta^2-34\omega^2=0,\qquad
 D_2=34cd+2\theta^2-34\omega^2=0.
\]
An explicit degree-two map $D\longrightarrow\Gamma$ is
\begin{equation}\label{eq:D-Gamma-map}
 (z_0,z_1,z_2,z_3)
 =(17c^2+d^2,\,4\theta\omega,\,d^2-17c^2,\,34cd).
\end{equation}
Direct substitution gives the two equations of $\Gamma$, and on $D$
it gives the twist identity
\begin{equation}\label{eq:xi-twist-exact}
 f_\Gamma\circ\pi_{D\Gamma}=-2(34c\theta)^2.
\end{equation}
For clarity, the first substituted equation differs from zero by
\[
 D_1(17c^2+d^2+2\theta^2+34\omega^2)
 -D_2(34cd-2\theta^2+34\omega^2),
\]
and the second substituted equation is identically zero.
The coordinate quadrics have no common projective zero on $D$;
their pullback degree therefore also verifies that the map has degree two.

The relevant tangent product on $D$ is
\[
 f_D=68(17c-d+2\theta)(c+d-2\omega).
\]
The second twist $\eta$ has the following coordinate vector in the integral basis of $L$:
\[
\begin{gathered}
(8279,3481,2702,-3882,423,-5959,1183,-1901,-2290,\\
 6568,-1452,-5100,-3098,-1659,3872,-2919,-5548,2753)^{\mathsf T}.
\end{gathered}
\]
It is totally positive and satisfies
\[
 N_{L/\mathbb Q}(\eta)=2^8\cdot17,
 \qquad (\eta)=\mathfrak p_{2,17}^{8}\mathfrak p_{17}.
\]
The relative discriminant of $L(\sqrt\eta)/L$ has norm $17$.
The exact twist-membership calculation also verifies the defining
affine Selmer conditions.

The second lift is the intersection
\[
 E:\quad
 E_1=34u^2+4uv+\eta\vartheta^2-68\eta\varpi^2=0,
 \qquad
 E_2=17u^2-4v^2-34\eta\vartheta^2+68\eta\vartheta\varpi=0.
\]
An explicit degree-two map $E\longrightarrow D$ is
\begin{equation}\label{eq:E-D-map}
\begin{aligned}
 c&=-68u^2-8uv,\\
 d&=-34u^2+8v^2,\\
 \theta&=-17u^2-68uv-4v^2,\\
 \omega&=\eta(\vartheta^2-68\vartheta\varpi+68\varpi^2).
\end{aligned}
\end{equation}
On $E$ the first two expressions also equal
\[
 c=\eta(2\vartheta^2-136\varpi^2),\qquad
 d=\eta(-68\vartheta^2+136\vartheta\varpi).
\]
These two parametrizations verify both equations of $D$ and yield
\begin{equation}\label{eq:eta-twist-exact}
 f_D\circ\pi_{ED}
 =\eta\bigl(136(17u+2v)(\vartheta-2\varpi)\bigr)^2.
\end{equation}
There is again no common projective zero of the coordinate quadrics:
the first two force $u=v=0$, and the equations of $E$ then force
$\vartheta=\varpi=0$.  Hence this is a degree-two morphism.
These identities hold in a rational polynomial ring with $\eta$ treated as an indeterminate.

We also identify the Jacobians including their twists.  For the
coefficient matrices of $D$,
\[
 \det(\lambda D_1+D_2)=34^2(\lambda^2-1)(\lambda^2-17).
\]
After removing the square factor, the determinant quartic is
$y^2=\lambda^4-18\lambda^2+17$, whose Jacobian is exactly $C$.
For the coefficient matrices of $E$ (minus twice the displayed forms),
\[
 \det(\lambda E_1^{\rm mat}+E_2^{\rm mat})
 =2^8\,17\eta^2(\lambda^2+34\lambda+17)
                     (\lambda^2-34\lambda+17).
\]
Dividing by the square $(16\eta)^2$ gives
\[
 y^2=17\lambda^4-19074\lambda^2+4913.
\]
The Jacobian formula for an even quartic
$y^2=a\lambda^4+b\lambda^2+c$ is
\[
 Y^2=X^3-2bX^2+(b^2-4ac)X.
\]
Here it gives
\[
 Y^2=X^3+38148X^2+363483392X.
\]
The substitutions $X=34^2T$, $Y=34^3U$ give
$U^2=T(T+16)(T+17)$, and $t=T+16$ gives exactly $F$.
The determinant-quartic identification of the Jacobian, including
its field of definition, is the standard one recalled in Fisher~\cite{Fisher4}, Section~3.

The covering chain has Jacobians
\[
 E\ (F)\ \longrightarrow\ D\ (C)\ \longrightarrow
 \Gamma\ (F)\ \longrightarrow\ C,
\]
with successive Jacobian isogenies $\phi,\psi,\phi$.
Thus its composite is a locally soluble $2\phi$-covering once local
solubility is verified below; it is the lift required in the definition
of $\Theta_3$.

\subsection{The norm pushout and its torsion point}

Let $M=L(s)$, $s^2=17$.  Certified maximal orders give
\[
 \operatorname{disc}(M)=\operatorname{disc}(L)^2.
\]
Consequently $M/L$ is unramified at every finite prime, as follows from the certified maximal orders of $M$. Put $\lambda=-17+4s$, $\delta=-\lambda$. At this pencil root the singular quadric is
\[
 4z^2+\eta\bigl((34-\lambda)\vartheta^2
             -68\vartheta\varpi+68\lambda\varpi^2\bigr)=0,
 \qquad z=v-\lambda u/2.
\]
An exact rational point on this cone is obtained from a verified
solution of
\begin{equation}\label{eq:final-norm-conic}
 x^2+\lambda y^2=\eta z_*^2 \qquad\text{over }M.
\end{equation}
The following construction supplies the required point on the auxiliary conic.

Choose the specified $r,h\in L$ with $\eta=r^2-17h^2$, and put
$k=2h(r+17h)$.  Since $\delta^2-34\delta+17=0$,
\[
 (r+h\delta)^2-\eta=k\delta.
\]
The certificate supplies an exact nonzero point $(X,Y,Z,W)$, with $W=0$,
on
\[
 X^2-\eta Y^2-kZ^2+17\eta kW^2=0.
\]
Set
\[
\begin{aligned}
 d_0&=Z^2-17\eta W^2,&v_0&=XZ-\eta YsW,\\
 v_1&=YZ-XsW,&w_0&=r+h\delta,\\
 x&=w_0v_0-\eta v_1,&y&=kd_0,&z_*&=v_0-w_0v_1.
\end{aligned}
\]
The exact identity
\[
 x^2-\delta y^2-\eta z_*^2
 =\delta k d_0\bigl(X^2-\eta Y^2-kZ^2+17\eta kW^2\bigr)
\]
proves \eqref{eq:final-norm-conic}. The displayed point and its substitution identities provide an exact algebraic witness for the norm equation.

With
\[
 U_0=(3s-17)/2,\qquad V_0=-(5s+13)/2,
\]
one has $U_0^2+\lambda V_0^2=-4(34-\lambda)$.  Define
\[
 U=U_0x+\lambda V_0y,\quad V=V_0x-U_0y,\quad
 Q=\left(0,\eta z_*,\frac{U+V/2}{34-\lambda},\frac V{68}\right).
\]
Then $Q$ lies on the singular cone.  Its actual tangent coefficients,
with the fixed normalization used at every place, are
\[
 T=\eta\begin{pmatrix}
 (34-8s)z_*\\
 4z_*\\
 \bigl((-17+3s)x+(-119+33s)y\bigr)/2\\
 (-51+17s)x+(-221+51s)y
 \end{pmatrix}^{\mathsf T}.
\]
Writing $T=T_0+sT_1$, the pushout quadratic form is
\begin{equation}\label{eq:final-norm-pushout}
 f_E(P)=N_{M/L}(T(P))=(T_0P)^2-17(T_1P)^2.
\end{equation}
The power-basis coefficients, field embeddings, norm point and quadric matrices are included in the supplementary data. Substitution in the defining equations verifies the norm point and tangent.

It remains to identify the torsion point selected by this tangent pair.
The conjugate roots $-17\pm4s$ both come from the $(u,v)$ block.
Over a splitting field the two associated cone projections are
degree-two involutions.  Their composite changes both diagonalized
coordinates of that block in sign; projectively this is
\[
 (u,v,\vartheta,\varpi)\longmapsto
                    (u,v,-\vartheta,-\varpi).
\]
This is precisely the deck involution of \eqref{eq:E-D-map}.
The product of the two tangent forms represents the difference of the
two cone divisor classes, as in Fisher's Lemma~7.1.  Its associated
two-torsion translation is therefore the kernel of the quotient
$F\longrightarrow C$, namely $(-1,0)$.  

The other two quotients of
$F$ have equations
\[
 y^2=x(x^2+30x+289),\qquad y^2=x(x^2-66x+1),
\]
and neither has the $j$-invariant of $C$; this also checks uniqueness
of the kernel identification.    Dividing the homogeneous
quadratic pushout by the square of a linear form makes it a degree-zero
rational function and leaves every local squareclass unchanged.

\subsection{Local solubility and certified local values}

The model $E$ has good reduction outside $2$ and $17$.  Indeed $\eta$
is a unit there, and the determinant pencil has unit discriminant.
At every such finite place its smooth genus-one reduction has a point
over the finite residue field, and a smooth point lifts.  This establishes local solubility at every good finite place.

At either real embedding, take $v=0$, $\varpi=1$ and
\[
 \vartheta=\frac{68+2\sqrt{2329}}{69}>2,
 \qquad u^2=2\eta\vartheta(\vartheta-2)>0.
\]
Both equations of $E$ then hold.  Complex places are immediate.

At the two dyadic places and the unique place above $17$, a direct
multivariable Hensel argument gives local points from the specified centres.  Normalize a centre $Z$ so that its minimum coordinate valuation
is zero, let $m$ be the minimum valuation of the two quadric residuals,
and let $d$ be the valuation of a nonsingular two-by-two Jacobian minor.
The quadric matrices are integral.  The inequality $m>2d$ gives an
actual zero $Z^*$ with coordinate correction valuation at least $h=m-d$,
keeping the other two coordinates fixed.  The data are listed in Table~\ref{tab:4}.
\begin{table}[htbp]
\centering
\caption{Hensel data for local points on the higher covering.}\label{tab:4}
\begin{tabular}{lrrrr}
\toprule
$\text{place}$ & $e$ & $m$ & $d$ & $h$\\
\midrule
$\mathfrak p_{2,17}$ & $17$ & $1929$ & $182$ & $1747$\\
$\mathfrak p_{2,1}$ & $1$ & $113$ & $10$ & $103$\\
$\mathfrak p_{17}$ & $18$ & $1171$ & $28$ & $1143$\\
\bottomrule
\end{tabular}
\end{table}
All inequalities are strict.  The positive correction precision also
preserves primitivity.  The strict Hensel inequalities therefore establish the existence of local points.

For the two dyadic places, let $c_f$ be the minimum valuation of the
coefficients of the fixed pushout matrix, and let $v_f=v(f_E(Z))$.
Since $Z$ is integral, the correction changes $f_E(Z)$ by valuation
at least $h+c_f$.  Consequently its squareclass is fixed whenever
$h+c_f-v_f>2e$.  

The local values are listed in Table~\ref{tab:5}.
\begin{table}[htbp]
\centering
\caption{Dyadic quadratic-pushout values and Hilbert-symbol factors.}\label{tab:5}
\begin{tabular}{lrrrr}
\toprule
$\text{place}$ & $c_f$ & $v_f$ & $h+c_f-v_f-2e$ & $(f_E(Z^*),b)$\\
\midrule
$\mathfrak p_{2,17}$ & $-8$ & $146$ & $1559$ & $+1$\\
$\mathfrak p_{2,1}$ & $0$ & $14$ & $87$ & $+1$\\
\bottomrule
\end{tabular}
\end{table}
The valuations in the table use primitive projective coordinates; projective scaling multiplies the quadratic pushout by a square.

At the degree-seventeen dyadic completion, $b$ is nonsquare.  Eighteen
independent exact norms $x_i^2-by_i^2$ span an eighteen-dimensional
subspace of its nineteen-dimensional squareclass space.  Hence they
span the full norm hyperplane, and its unique nonzero annihilator
is the Hilbert character $(\,\cdot\,,b)$.  This character gives the first sign in the table.  At the degree-one
completion the valuation and unit residue modulo eight give the
second sign by the usual quadratic Hilbert-symbol formula.
At the place above $17$, $b$ is a verified local square.  At both real
places $b>0$.  These places, and all complex places, therefore contribute
$+1$ once local solubility has been established.

\subsection{All good odd places, including coefficient content}

To compute the contributions at good odd places, we retain the coefficient content of the chosen global tangent.
Let
\[
 \mathfrak C_T=(T_1,T_2,T_3,T_4)
\]
be the fractional coefficient ideal of the actual tangent in the
original $E$ coordinates, as an ideal of $M$.  At a prime
$\mathfrak q\nmid2\cdot17$, the singular quadric is a good rank-three
cone, and its vertex does not meet the reduction of $E$; otherwise
the intersection of quadrics would be singular.  Projection from the
vertex therefore takes a primitive local point to a primitive point
on a smooth conic.  A primitive tangent to that conic has an evaluation of even valuation.  

Restoring the coefficient content yields
\[
 v_{\mathfrak q}(T(P))
       \equiv v_{\mathfrak q}(\mathfrak C_T)\pmod2.
\]
The target $b$ has even valuation at all these places.  Hilbert-symbol
projection under the norm now gives the complete good odd-place product
\begin{equation}\label{eq:odd-content-product}
 \prod_{\mathfrak q\nmid2\cdot17}
 (\pi_{\mathfrak q},b)_{M_{\mathfrak q}}^
                              {v_{\mathfrak q}(\mathfrak C_T)}.
\end{equation}
At an inert prime the relative residue degree is two and the unit
residue of $b$ becomes a square, so the contribution is $+1$.
At split primes the finite-field quadratic character determines it.

Table~\ref{tab:6} groups the complete factorization of
$\mathfrak C_T$ by rational prime.  Each listed pair is
$(\text{residue degree},\text{content valuation})$; repeated pairs
refer to distinct prime ideals, whose complete records are in the
certificate.
\begin{table}[htbp]
\centering
\caption{Factorization of the tangent coefficient ideal by rational prime.}\label{tab:6}
\begin{tabular}{lr}
\toprule
$\text{rational prime}$ & $(f,v_{\mathfrak q}(\mathfrak C_T))$\\
\midrule
$2$ & $(1,-2),(1,11),(1,1)$\\
$17$ & $(1,6),(1,6)$\\
$113$ & $(2,-28),(16,-22),(16,-22)$\\
$149$ & $(4,-49),(4,-49)$\\
$7159$ & $(1,-71),(1,1),(1,-71)$\\
$10228457$ & $(1,-10),(1,-50),(1,-10),(1,-50)$\\
$33062332448349229$ & $(1,1)$\\
\bottomrule
\end{tabular}
\end{table}
Every rational prime is proved prime.  Multiplication of the eighteen
prime-ideal powers reconstructs the actual coefficient ideal, and each
exponent also equals the minimum of the four tangent-coefficient
valuations.  These identities certify the complete support and exponents in the table.
Only six good odd entries have odd exponent; their factors are listed in Table~\ref{tab:7}.
\begin{table}[htbp]
\centering
\caption{Good odd-place factors with odd coefficient-content valuation.}\label{tab:7}
\begin{tabular}{lrrr}
\toprule
$p$ & $f$ & $v_{\mathfrak q}(\mathfrak C_T)$ & $(\pi_{\mathfrak q},b)_{M_{\mathfrak q}}$\\
\midrule
$149$ & $4$ & $-49$ & $-1$\\
$149$ & $4$ & $-49$ & $-1$\\
$7159$ & $1$ & $-71$ & $+1$\\
$7159$ & $1$ & $1$ & $-1$\\
$7159$ & $1$ & $-71$ & $+1$\\
$33062332448349229$ & $1$ & $1$ & $+1$\\
\bottomrule
\end{tabular}
\end{table}
Their product is $-1$.  In particular the two factors above $149$
cancel, and the net negative factor lies above $7159$.
The two relevant degree-one primes of $L$ above $7159$ have residues
$\alpha=1338,6724$.  Above $1338$ the two $s$-residues are
$4547,2612$, with content valuations $0,1$ and target character $-1$;
above $6724$ the conjugate content valuations are both $-71$ and
the target character is $+1$.

The coefficient ideal is the Hermite normal form of the $36$-by-$144$ matrix formed from the four multiplication-by-coefficient maps. Residue-field quadratic characters give the signs displayed above. These calculations use the maximal order of $M$.
The content calculation accounts for the local normalization of the tangent.

\subsection{The global pairing value}

Let $e\in\operatorname{Sha}(F/L)$ be the class represented by $E$.  The explicit
locally soluble chain gives $2e=a$ in $\operatorname{Sha}(F/L)$, and $E\longrightarrow C$
is the required $2\phi$-covering.  Since $b\in S_3$, its image in
$\operatorname{Sha}(F/L)$ is also of the form $2b_1$ for a suitable lifted class.
Fisher's definition gives
\[
 \Theta_3(a,b)=\langle e,b\rangle
             =\langle a,b_1\rangle.
\]
If another locally soluble lift $e'$ is chosen, then $2(e-e')=0$;
adjointness of multiplication by two under the Cassels--Tate pairing
therefore yields
\[
 \langle e-e',b\rangle
 =\langle e-e',2b_1\rangle
 =\langle2(e-e'),b_1\rangle=0.
\]
This proves independence of all lift choices at once.  Equivalently,
the first-lift ambiguity lies in $S'_1=Q'$ and has trivial Sha image;
after fixing the first lift, the second-lift ambiguity lies in
$\operatorname{Sha}(F/L)[\phi]$, and the condition $b\in\ker\Theta_1$ annihilates it.
This accounts for both stages of the lift.

The pushout selects $F[\phi]$ as proved above, so Fisher's Hilbert-symbol
formula applies to \eqref{eq:final-norm-pushout} and the same target $b$.
The contributions are listed in Table~\ref{tab:8}.
\begin{table}[htbp]
\centering
\caption{Local contributions to the higher Cassels--Tate pairing.}\label{tab:8}
\begin{tabular}{lr}
\toprule
$\text{places}$ & $\text{product}$\\
\midrule
$\mathfrak p_{2,17}$ & $+1$\\
$\mathfrak p_{2,1}$ & $+1$\\
$\text{all good odd places}$ & $-1$\\
$\mathfrak p_{17}$ & $+1$\\
$\text{archimedean places}$ & $+1$\\
\bottomrule
\end{tabular}
\end{table}
Hence
\[
 \Theta_3(-17,b)=(+1)(+1)(-1)(+1)(+1)=-1,
\]
which proves \eqref{eq:theta3-nonzero} and therefore
\eqref{eq:higher-rank-three}.

The finite Selmer groups, explicit locally soluble coverings and the Cassels--Tate image criterion in Fisher's Theorem 3.1 give the rank bound. The theorem places no finiteness requirement on $\operatorname{Sha}$, and applies here to $\operatorname{Sha}(F/L)$ and $\operatorname{Sha}(C/L)$.

% END higher_pairing_section.tex


% BEGIN chabauty_section.tex
\section{Three points and local analysis at $5$}\label{sec:points}

\subsection{Exact points and the torsion subgroup}

Finite-index questions in the computation of Mordell--Weil groups are treated by Siksek \cite{Siksek}. Here the reduction tests establish the saturation statements required for the local closure.

Over the field $\Q[\tau]/(F)$ of Section~\ref{sec:field}, the supplementary data specify points $P_1,P_2,P_3$ on $A$ by their power-basis coordinates. Substitution in the curve equation and the modular map gives
\begin{equation}\label{eq:threej}
j(P_1)=\frac{9261}{8},\qquad
j(P_2)=-\frac{132651}{2},\qquad
j(P_3)=-\frac{1167051}{512}.
\end{equation}
All three values are $2$ modulo $5$; their residues modulo $25$ are respectively $17,12,2$. These congruences place the three points in the residue class imposed by the Frey curve.

The polynomial $F$ has simple roots modulo $11$ and modulo $37$.
Thus $L$ has good reduction places with residue fields $\F_{11}$ and
$\F_{37}$ for $A$. Exact counting gives
\[
\#A(\F_{11})=12,\qquad \#A(\F_{37})=40.
\]
Prime-to-residue-characteristic torsion injects under good reduction.
Using both places, including for the $11$- and $37$-primary parts, shows
that $\#A(L)_{\rm tors}$ divides $4$. The point $T=(7,13)$ satisfies
\[
2T=(11/4,-15/8),\qquad4T=O,
\]
and $2T\ne O$. Consequently
\begin{equation}\label{eq:tors}
A(L)_{\rm tors}=\langle T\rangle\simeq\Z/4\Z.
\end{equation}

Put $H=\langle P_1,P_2,P_3,T\rangle$. At places above the rational primes $11,37,97$, each of the $15$ nonzero binary combinations of $P_1,P_2,P_3,T$ has a reduction that is not divisible by $2$ at at least one selected place. Hence these four classes are independent
in $A(L)/2A(L)$.

These independent reduction classes prove that $H$ is $2$-saturated. Indeed, if $2Q=a_1P_1+a_2P_2+a_3P_3+kT$, the four coefficients must be even. Subtracting the corresponding half in $H$ leaves an element of $A(L)[2]=\{O,2T\}\subset H$. Thus $Q\in H$.

\subsection{The exact finite group at $5$}

Since $F$ is irreducible and separable modulo $5$, the field
$K=L\otimes_{\Q}\Q_5$ is the unramified extension of degree $18$.
Its integer ring is $\OO=\Z_5[\tau]$ and its residue field is
$k=\F_{5^{18}}$. In particular $1,\tau,\ldots,\tau^{17}$ is an
integral basis over $\Z_5$.

The curve $A$ has good reduction at $5$ and $a_5(A)=-2$. The usual
Frobenius recurrence, also checked by finite-field group arithmetic, gives
\begin{align*}
\#A(k)&=3814695421472\\
&=2^5\cdot13\cdot37\cdot73\cdot181\cdot18757.
\end{align*}
The exact order of $\overline P_1$ is
\[
m=953673855368=\#A(k)/4.
\]
For each prime divisor $\ell$ of $m$, the certificate verifies
$[m/\ell]\overline P_1\ne O$, as well as $[m]\overline P_1=O$.
Furthermore,
\[
\overline T\notin\langle\overline P_1\rangle,
\qquad 2\overline T=[m/2]\overline P_1.
\]
Thus $\langle\overline P_1,\overline T\rangle$ has index $2$ in $A(k)$.

The fibre $j=2$ over $k$ consists of exactly the $18$ Frobenius conjugates
of $\overline P_1$. Its degree-$18$ fibre polynomial is irreducible
modulo $5$. The Pohlig--Hellman computation expresses all $18$ points in $\langle\overline P_1,\overline T\rangle$. In particular
\[
\overline P_2=[n_2]\overline P_1+\overline T,
\qquad n_2=18179330055,
\]
\[
\overline P_3=[n_3]\overline P_1+\overline T,
\qquad n_3=237047023703.
\]
The Frobenius indices of $\overline P_2$ and $\overline P_3$ are $11$
and $7$. The supplementary data give the expressions and order witnesses for all eighteen points.

\subsection{Three primitive formal directions}

For the formal-group and local-field conventions used in this calculation, see Silverman \cite{SilvermanAdvanced,Silverman}.

Define
\[
R_1=[m]P_1,\qquad R_2=P_2-[n_2]P_1-T,\qquad
R_3=P_3-[n_3]P_1-T.
\]
They lie in the formal group $A_1(\OO)$, the kernel of reduction.
Use the short model
\[
u=4x-1,\quad v=8y+4x+4,\qquad v^2=u^3-11u-890,
\]
and the formal parameter $z=-u/v$. If $\Lambda_i=\log(R_i)$, then
$\Lambda_i\in5\OO$ and $\log z\equiv z\pmod{25}$ on $5\OO$. Here the logarithm denotes the formal-group power series.
Exact calculation of $z(R_i)$ modulo $25$ gives an $18\times3$ matrix
for $\Lambda_i/5$ modulo $5$. Its rows of coefficient indices $0,1,2$
are
\begin{equation}\label{eq:logminor}
\begin{pmatrix}3&1&2\\4&2&3\\0&1&2\end{pmatrix},
\qquad\det=3\quad\text{in }\F_5.
\end{equation}
Therefore the three logarithms are independent over $\Q_5$, and the
formal lattice they generate is primitive in $5\OO$ (equivalently,
the $\Lambda_i/5$ span a primitive submodule of $\OO$). Since
\[
\log(mP_1)=\Lambda_1,
\quad\log(mP_2)=n_2\Lambda_1+m\Lambda_2,
\quad\log(mP_3)=n_3\Lambda_1+m\Lambda_3,
\]
and $m$ is a $5$-adic unit, the $P_i$ are independent in $A(L)$.
Combining this lower bound with the preceding descent proves
\begin{equation}\label{eq:rank3}
\rank A(L)=3.
\end{equation}

The same logarithm calculation also proves global $5$-saturation. Suppose
$5Q=a_1P_1+a_2P_2+a_3P_3+kT$ with $Q\in A(L)$.
Multiplication by $m$ kills $T$. Since $5\nmid\#A(k)$, reduction of
$5mQ$ to zero implies that $mQ\in A_1(\OO)$. Taking formal logarithms,
dividing by $5$, and reducing modulo $5$ forces each $a_i$ to be
divisible by $5$, by \eqref{eq:logminor} and the unit change of basis
above. Also $5T=T$. Subtracting an element of $H$ therefore leaves a
$5$-torsion point, which is zero by \eqref{eq:tors}. Hence $Q\in H$.

The two saturation results and the rank computation show that the subgroup generated by the three points and the torsion has finite index prime to ten. This index condition identifies its local closure with that of the full Mordell--Weil group in the next subsection.

\subsection{The local closure}

By \eqref{eq:rank3}, $d=[A(L):H]$ is finite. The two saturation arguments give $\gcd(d,10)=1$; the closure argument uses this condition regardless of whether $d=1$. Let
\[
G_e=A(\OO/5^e\OO),\quad H_e=\operatorname{im}(H\to G_e),
\quad M_e=\operatorname{im}(A(L)\to G_e).
\]
Properness and good reduction make these reduction maps valid for every
point, irrespective of its affine-coordinate denominators. The kernel of
$G_e\to G_1$ is a $5$-group. Since $[G_1:H_1]=2$, the index
$[G_e:H_e]$ has the form $2\cdot5^{c_e}$. On the other hand,
$M_e/H_e$ is a quotient of $A(L)/H$. Thus
\[
|M_e/H_e|\mid d\quad\text{and}\quad |M_e/H_e|\mid2\cdot5^{c_e},
\]
which forces $M_e=H_e$ for every $e$. Consequently $A(L)$ and $H$
have exactly the same closure in $A(K)$.

The kernel $H\cap A_1(\OO)$ is generated by
\[
D=[m/2]P_1+2T,\qquad R_2,\qquad R_3,
\qquad 2D=R_1.
\]
To see this, subtract the appropriate multiples of $R_2,R_3$ from an
element of $H$. What remains is $aP_1+kT$. Its reduction is zero
precisely when $k$ is even and the remaining relation is generated by
$m\overline P_1=O$ and $2\overline T=[m/2]\overline P_1$.
Since $2$ is a unit in $\Z_5$, the closure of this formal kernel is
therefore parametrized exactly by the three logarithms $\Lambda_i$.

For each of the eighteen relevant residue points choose the specified
base point $Q_i=[n_i]P_1+[k_i]T$. Every global point in that residue
disk belongs to the parametrized set
\begin{equation}\label{eq:param}
Q_i(\mathbf a)=Q_i+\exp(a_1\Lambda_1+a_2\Lambda_2+a_3\Lambda_3),
\qquad\mathbf a\in\Z_5^3.
\end{equation}
Here exp is the inverse of the formal-group logarithm on the formal neighborhood of the origin. Independence of the three logarithms makes the parameter map injective.

\subsection{The exact calculation modulo $25$}

Write $j(Q_i(\mathbf a))$ in the integral basis
$1,\tau,\ldots,\tau^{17}$. Its nonconstant coefficients vanish
modulo $5$ on every relevant disk. Divide them by $5$ and denote the
seventeen resulting functions by $F_{i,1},\ldots,F_{i,17}$.
If $j(Q_i(\mathbf a))\in\Q_5$, all seventeen must vanish.

Modulo $5$ they have the affine-linear form
\begin{equation}\label{eq:mod25system}
F_{i,r}(\mathbf a)\equiv c_{i,r}
 +B_{i,r1}a_1+B_{i,r2}a_2+B_{i,r3}a_3\pmod5.
\end{equation}
These coefficients are calculated by exact arithmetic in $\OO/25\OO$.
On the short model a formal displacement $z\in5\OO/25\OO$ gives
\[
u(Q+R)=u(Q)+2v(Q)z,\qquad
v(Q+R)=v(Q)+(3u(Q)^2-11)z.
\]
The coordinate changes follow by expanding the affine addition formulas modulo $25$.

Every $17\times3$ matrix $B_i$ has rank $3$. For fifteen disks the
seventeen equations are inconsistent. For the other three the complete result is listed in Table~\ref{tab:9}
\begin{table}[htbp]
\centering\caption{The three residue disks surviving the calculation modulo 25.}\label{tab:9}
\begin{tabular}{cccr}
\toprule
Frobenius index $i$&Unique parameter modulo $5$&Exact point&$j\bmod25$\\\midrule
0&$(0,0,0)$&$P_1$&17\\
7&$(0,0,1)$&$P_3$&2\\
11&$(0,1,0)$&$P_2$&12\\\bottomrule
\end{tabular}
\end{table}
For each excluded disk, Appendix~\ref{app:disks} gives three rows with nonzero determinant, the unique parameter solving them, and a nonzero residual in one further equation. The supplementary data contain the full matrices.

\subsection{Analytic uniqueness on each residue disk}

The finite calculation extends to each entire residue disk through the integral Taylor expansion of the modular invariant. On the eighteen relevant disks, $x-7$ is a unit, so \eqref{eq:jmap} is an integral regular function. The translation-invariant derivation on the short model is
\[
\mathcal D=2v\frac\partial{\partial u}
 +(3u^2-11)\frac\partial{\partial v}.
\]
It preserves the integral coordinate ring localized at the unit
denominator of $j$. Taylor expansion in the formal logarithm $\lambda$
therefore has coefficients $\mathcal D^n j/n!$ with integral numerators.
Here $\lambda=\sum a_i\Lambda_i$ and every $\Lambda_i\in5\OO$.


Thus a term of total parameter degree $n$ has valuation at least
\[
n-v_5(n!).
\]
For $n\ge2$ this is at least $2$ and tends to infinity with $n$.
After dividing a nonconstant coefficient by $5$, all nonlinear terms
belong to $5\Z_5$. The resulting functions are restricted power series
on the whole of $\Z_5^3$, with constant derivative modulo $5$ equal
to the matrix in \eqref{eq:mod25system}.

Choose the three rows giving an invertible minor at a surviving disk.
Their vector of functions has the form
\[
\mathbf F(\mathbf a)=\mathbf c+B\mathbf a+5\mathbf G(\mathbf a),
\]
where $B\in M_3(\Z_5)$ is invertible modulo $5$, and $\mathbf G$
is an integral restricted power series, hence is $1$-Lipschitz.
If $\mathbf F(\mathbf a)=\mathbf F(\mathbf b)=0$, reduction modulo
$5$ gives $\mathbf a\equiv\mathbf b$. If they differ, put
$\nu=\min_jv_5(a_j-b_j)$ and write
$\mathbf a-\mathbf b=5^\nu\mathbf h$ with $\mathbf h\not\equiv0$.
Then
\[
0=\mathbf F(\mathbf a)-\mathbf F(\mathbf b)
 \equiv5^\nu B\mathbf h\pmod{5^{\nu+1}},
\]
which is impossible. Each surviving disk therefore contains at most one
parameter whose $j$-value is in $\Q_5$. The corresponding exact point
$P_1,P_3,P_2$ supplies such a parameter, so it is the only one.

This proves uniqueness on the entire parameter space $\Z_5^3$.

\subsection{Exceptional points and pole disks}

The fibre calculation checks that its degree-$18$ polynomial is coprime
to $b(x)$ and $x-7$ modulo $5$. Thus each of its roots supplies exactly
one point in the relevant affine chart.
At $O$, the function $j$ has a pole of order $1$ with unit leading
coefficient. At $T$ it has a pole of order $17$ with unit leading
coefficient. Points strictly inside these pole disks have negative
$5$-adic valuation of $j$ and cannot satisfy \eqref{eq:jmod5}.

At $-T=(7,-21)$ the apparent $0/0$ in \eqref{eq:jmap} is removable.
The exact norm elimination identity contains $(x-7)^{17}$, while the
numerator on the other branch is a unit. It follows that the numerator
is divisible by $(x-7)^{17}$ in the integral local power-series ring on
the $-T$ branch. Its cancelled constant term is $3$ modulo $5$.
The entire $-T$ disk is therefore irrelevant. These checks, together
with properness of $A$, exhaust the possible reductions of a global
point.

\begin{proposition}\label{prop:rationalj}
For the specified embedding $L\hookrightarrow K$, the set
\[
\{Q\in A(L):j(Q)\in\Q_5,\ j(Q)\equiv2\pmod5\}
\]
is exactly $\{P_1,P_2,P_3\}$. In particular the rational $j$-values
in this residue class are precisely the three values in
\eqref{eq:threej}.
\end{proposition}

\section{Primitive solutions and rational points}\label{sec:applications}

The modular-point classification completes the proof of Theorem 1.1 as follows.

\subsection{Rational parameters and the power condition}

Put
\[
s=\frac{X^3}{Z^{17}},\qquad s-2=\left(\frac YZ\right)^{17}.
\]
The nonzero assumptions imply $s\ne0,2$, and the modular invariant is
\[
J(s)=\frac{27s(9s-16)^3}{8(s-2)}.
\]
For each candidate $j$, solve $27s(9s-16)^3-8j(s-2)=0$. The three exact factorizations are
\[
\begin{array}{l}
27s(9s-16)^3-9261(s-2)\\
\qquad=27(s-1)(9s-2)(81s^2-333s+343),
\end{array}
\]
\[
\begin{array}{l}
27s(9s-16)^3+530604(s-2)\\
\qquad=27(s+2)(729s^3-5346s^2+17604s-19652),
\end{array}
\]
\[
\begin{array}{l}
27s(9s-16)^3+\frac{1167051}{64}(s-2)\\
\qquad=\frac3{64}(4s+1)(104976s^3-586116s^2+1141857s-778034).
\end{array}
\]
The quadratic factor has discriminant $-243$. The cubic factors reduce modulo $7$ to
\[
s^3+2s^2+6s+4,\qquad4s^3+s^2+3s+2,
\]
respectively. Neither has a root in $\mathbb F_7$, so both cubic factors are irreducible over $\mathbb Q$. It follows that
\[
s\in\left\{1,\frac29,-2,-\frac14\right\}.
\]
For the last three values, the corresponding values of $s-2$ are
\[
-\frac{16}{9},\qquad-4,\qquad-\frac94.
\]
Their $2$-adic valuations are $4$, $2$ and $-2$. None is divisible by $17$, so none is a seventeenth power in $\mathbb Q$. Therefore $s=1$, and $(Y/Z)^{17}=-1$ gives $Y=-Z$. Primitivity forces $|Y|=|Z|=1$, and the original equation gives exactly the two stated triples. This proves Theorem~\ref{thm:ternary}.

\subsection{Consequences and affine rational points}

The primitive classification gives the following consequences.

\begin{corollary}\label{cor:fixedN}
For every positive integer $N$, the equation $X^3-Y^{17}=2N^{17}$ has no solution in positive integers $X,Y$ with $\gcd(Y,N)=1$.
\end{corollary}
\begin{proof}
Take $Z=N$ in Theorem~\ref{thm:ternary}.
\end{proof}

\begin{corollary}\label{cor:plus}
The only solution of $X^3+Y^{17}=2Z^{17}$ in positive integers with $\gcd(Y,Z)=1$ is $(X,Y,Z)=(1,1,1)$.
\end{corollary}
\begin{proof}
Apply Theorem~\ref{thm:ternary} to the triple $(X,-Y,Z)$.
\end{proof}

\begin{corollary}\label{cor:other}
The equation $X^3+2Y^{17}=Z^{17}$ has no solution in positive integers with $\gcd(Y,Z)=1$.
\end{corollary}
\begin{proof}
Apply Theorem~\ref{thm:ternary} to $(X,Z,-Y)$.
\end{proof}

\begin{corollary}\label{cor:rational}
The affine curve $x^3-y^{17}=2$ has exactly one rational point:
\[
\{(x,y)\in\mathbb Q^2:x^3-y^{17}=2\}=\{(1,-1)\}.
\]
\end{corollary}
\begin{proof}
Neither coordinate can vanish, since neither $2$ nor $-2$ is the required rational perfect power. For a prime $\ell$ at which a coordinate has negative valuation, the equation forces both terms to have equal negative valuation, so $3v_\ell(x)=17v_\ell(y)$. Coprimality of $3$ and $17$ therefore gives integers $A,B$ and a positive integer $T$ such that
\[
x=\frac A{T^{17}},\qquad y=\frac B{T^3},\qquad\gcd(B,T)=1.
\]
Clearing denominators gives
\[
A^3=B^{17}+2(T^3)^{17}.
\]
Theorem~\ref{thm:ternary}, applied to $(A,B,T^3)$, gives $T=1$, $A=1$ and $B=-1$.
\end{proof}


\appendix
\setcounter{table}{0}\renewcommand{\thetable}{A.\arabic{table}}\renewcommand{\theHtable}{A.\arabic{table}}
\section{Exact data displayed in the text}

% BEGIN generated_tables.tex
\subsection{Coefficients of the modular map}\label{app:jmap}
Write $a(x)=\sum a_i x^i$ and $b(x)=\sum b_i x^i$ in \eqref{eq:jmap}. The coefficients are listed in Table~\ref{tab:A-1}:
\begingroup\small
\begin{longtable}{r r r}
\caption{Coefficients of the polynomials in the modular j-map.}\label{tab:A-1}\\
\toprule $i$&$a_i$&$b_i$\\\midrule\endfirsthead
\toprule $i$&$a_i$&$b_i$\\\midrule\endhead
0&16619703086561043840&393689605817603264\\
1&2560483516165554624&-7141870055640467776\\
2&-27272162241361493664&-3232743977333766672\\
3&-14792936512811763984&3019362948630772592\\
4&4121460878619558792&2856893370701322548\\
5&7590065519042974644&361912536250488708\\
6&2206827523498348226&-503002839579031251\\
7&-629465475103177139&-236871273757462563\\
8&-591397116666291076&-14302404833132003\\
9&-114333139654357588&17129280335934548\\
10&20194867810083620&5158593525523346\\
11&13102383326408750&531227490780958\\
12&2098022038690768&20678454031406\\
13&128196606784368&272231243744\\
14&2883334192042&862568081\\
15&19441608861&196765\\
16&21602108&1\\
17&744&0\\
\bottomrule\end{longtable}\endgroup
\subsection{Finite exclusion witnesses for the eighteen disks}\label{app:disks}
All entries in Table~\ref{tab:A-2} are modulo $5$. The rows column gives coefficient indices for a nonsingular minor. The parameter is the unique solution of those three equations. An excluded disk has the displayed nonzero residual in a further equation. The supplementary data contain the full coefficients of that equation.
\begingroup\small
\begin{longtable}{r c r c c r}
\caption{Linear-system witnesses modulo five for the 18 residue disks.}\label{tab:A-2}\\
\toprule Disk&Rows&Det.&Parameter&Extra row&Residual\\\midrule\endfirsthead
\toprule Disk&Rows&Det.&Parameter&Extra row&Residual\\\midrule\endhead
0&$(1,2,3)$&1&$(0,0,0)$&--&--\\
1&$(1,2,3)$&4&$(1,1,4)$&4&4\\
2&$(1,2,3)$&2&$(0,3,3)$&5&1\\
3&$(1,2,3)$&3&$(3,0,2)$&4&2\\
4&$(1,2,4)$&4&$(0,2,4)$&3&3\\
5&$(1,2,3)$&2&$(1,0,4)$&4&2\\
6&$(1,2,3)$&2&$(0,3,4)$&4&3\\
7&$(1,3,5)$&2&$(0,0,1)$&--&--\\
8&$(1,2,3)$&2&$(0,1,3)$&4&4\\
9&$(1,2,3)$&1&$(4,3,1)$&4&1\\
10&$(1,2,4)$&2&$(3,3,4)$&3&1\\
11&$(1,2,4)$&1&$(0,1,0)$&--&--\\
12&$(1,2,3)$&2&$(4,0,3)$&6&2\\
13&$(1,2,3)$&3&$(3,3,3)$&4&1\\
14&$(1,2,3)$&2&$(2,1,2)$&4&2\\
15&$(1,2,3)$&1&$(2,2,0)$&4&2\\
16&$(1,2,3)$&2&$(1,1,1)$&4&2\\
17&$(1,2,3)$&3&$(1,2,1)$&4&3\\
\bottomrule\end{longtable}\endgroup

% END generated_tables.tex


\section*{Supplementary data}
Online Resource 1 contains the exact field and ideal data, local point coordinates, descent matrices and residue-disk systems used in the proof. Number-field calculations use PARI/GP \cite{PARI}. \file{Supplementary_Data_Guide.txt} identifies the supporting files for each section and gives the software requirements and reproduction commands.

For the class group, \file{active_L18/class_unit_2part_standalone_verify.gp} checks the unit and ideal identities, and \file{active_L18/verify_factor_base_fixed.gp} checks membership of the initial factor base. The expanded factor base is stored in \file{active_L18/class_certification/factor_base_50000.bin.} The twenty interval records and the analytic-bound data are in \file{active_L18/final_proof/reproduction_evidence/full_class/.}

The local Selmer matrices and lower pairing data are recorded in \file{active_L18/final_proof/proof_data.json} and the conic files in \file{active_L18/simon_descent/.} The higher pairing uses \file{active_L18/simon_descent/Cassels_E_final_exact_data.json} and \file{verify_Cassels_higher_final.gp}. Its covering maps and Hensel inequalities have separate checking programs identified in the data guide.

The points and modular map are in \file{active_L18/final_proof/points_and_jmap.json.} The finite group and the complete systems modulo twenty-five are in \file{L18_finite_group5_certificate.json}, \file{L18_mod25_dimension3_certificate.json} and \file{L18_mod25_dimension3_compact_certificate.json}, all under \file{active_L18/.} The portable verifier checks the printed finite identities and matrices; the class-generation and local-image calculations have their own programs. The analytic uniqueness argument is given in Section 7.6. The additional finite calculations for Section 8 are supplied in the \file{ternary_application} directory of Online Resource 1.

\section*{Statements and Declarations}
\subsection*{Data and code availability}
The exact computational data and programs supporting this article are provided in Online Resource 1. \file{Supplementary_Data_Guide.txt} gives the section-by-section file index, software requirements and reproduction commands.

\subsection*{Funding}
This research did not receive any specific grant from funding agencies in the public, commercial, or not-for-profit sectors.

\subsection*{Competing interests}
The author declares no known competing financial interests or personal relationships that could have appeared to influence the work reported in this paper.

\subsection*{Author contributions}
The sole author is responsible for the research and manuscript.

\subsection*{AI-assisted editing and computational checks}
Codex was used for language editing, manuscript formatting, and checks of selected computational steps and verification programs. The author is responsible for the mathematical arguments, computations, and final manuscript.

\begin{thebibliography}{99}

\bibitem{BDDF}
Belabas, K., Diaz y Diaz, F., Friedman, E.: \emph{Small generators of the ideal class group}. Math. Comp. 77(262), 1185--1197 (2008). \url{https://doi.org/10.1090/S0025-5718-07-02003-0}.

\bibitem{BVY}
Bennett, M.A., Vatsal, V., Yazdani, S.: \emph{Ternary Diophantine equations of signature $(p,p,3)$}. Compos. Math. 140(6), 1399--1416 (2004). \url{https://doi.org/10.1112/S0010437X04000983}.

\bibitem{BCDT}
Breuil, C., Conrad, B., Diamond, F., Taylor, R.: \emph{On the modularity of elliptic curves over Q: wild 3-adic exercises}. J. Amer. Math. Soc. 14(4), 843--939 (2001). \url{https://doi.org/10.1090/S0894-0347-01-00370-8}.

\bibitem{BLS}
Br\"oker, R., Lauter, K., Sutherland, A.V.: \emph{Modular polynomials via isogeny volcanoes}. Math. Comp. 81(278), 1201--1231 (2012). \url{https://doi.org/10.1090/S0025-5718-2011-02508-1}.

\bibitem{Bruin}
Bruin, N.: \emph{Chabauty methods using elliptic curves}. J. Reine Angew. Math. 562, 27--49 (2003). \url{https://doi.org/10.1515/crll.2003.076}.

\bibitem{CasselsIII}
Cassels, J.W.S.: \emph{Arithmetic on curves of genus 1. III. The Tate--\v{S}afarevi\v{c} and Selmer groups}. Proc. London Math. Soc. (3) 12(1), 259--296 (1962). \url{https://doi.org/10.1112/plms/s3-12.1.259}.

\bibitem{CasselsIV}
Cassels, J.W.S.: \emph{Arithmetic on curves of genus 1. IV. Proof of the Hauptvermutung}. J. Reine Angew. Math. 211, 95--112 (1962). \url{https://doi.org/10.1515/crll.1962.211.95}.

\bibitem{Cohen1993}
Cohen, H.: \emph{A Course in Computational Algebraic Number Theory}. Graduate Texts in Mathematics, vol. 138. Springer, Berlin (1993). \url{https://doi.org/10.1007/978-3-662-02945-9}.

\bibitem{Cohen2000}
Cohen, H.: \emph{Advanced Topics in Computational Number Theory}. Graduate Texts in Mathematics, vol. 193. Springer, New York (2000). \url{https://doi.org/10.1007/978-1-4419-8489-0}.

\bibitem{Coleman}
Coleman, R.F.: \emph{Effective Chabauty}. Duke Math. J. 52(3), 765--770 (1985). \url{https://doi.org/10.1215/S0012-7094-85-05240-8}.

\bibitem{DS}
Diamond, F., Shurman, J.: \emph{A First Course in Modular Forms}. Graduate Texts in Mathematics, vol. 228. Springer, New York (2005). \url{https://doi.org/10.1007/978-0-387-27226-9}.

\bibitem{Fisher4}
Fisher, T.: \emph{Some improvements to 4-descent on an elliptic curve}. In: van der Poorten, A.J., Stein, A. (eds.) Algorithmic Number Theory, Lecture Notes in Computer Science, vol. 5011, pp. 125--138. Springer, Berlin (2008). \url{https://doi.org/10.1007/978-3-540-79456-1_8}.

\bibitem{Fisher}
Fisher, T.: \emph{Higher descents on an elliptic curve with a rational 2-torsion point}. Math. Comp. 86(307), 2493--2518 (2017). \url{https://doi.org/10.1090/mcom/3163}.

\bibitem{FNS}
Freitas, N., Naskr\k{e}cki, B., Stoll, M.: \emph{The generalized Fermat equation with exponents 2, 3, n}. Compos. Math. 156(1), 77--113 (2020). \url{https://doi.org/10.1112/S0010437X19007693}.

\bibitem{JhaSahoo}
Jha, S., Sahoo, S.: \emph{Asymptotic Fermat equation of signature $(r,r,p)$ over totally real fields}. Ramanujan J. 68, Article 30 (2025). \url{https://doi.org/10.1007/s11139-025-01168-6}.

\bibitem{Kraus}
Kraus, A.: \emph{On the equation $x^p+y^q=z^r$: a survey}. Ramanujan J. 3, 315--333 (1999). \url{https://doi.org/10.1023/A:1009835521324}.

\bibitem{PARI}
The PARI Group: \emph{PARI/GP, version 2.19.0}. Software and documentation. Accessed 13 November 2026. \url{https://pari.math.u-bordeaux.fr/}.

\bibitem{Ribet}
Ribet, K.A.: \emph{On modular representations of $\operatorname{Gal}(\overline{\mathbb Q}/\mathbb Q)$ arising from modular forms}. Invent. Math. 100, 431--476 (1990). \url{https://doi.org/10.1007/BF01231195}.

\bibitem{SchaeferStoll}
Schaefer, E.F., Stoll, M.: \emph{How to do a $p$-descent on an elliptic curve}. Trans. Amer. Math. Soc. 356(3), 1209--1231 (2004). \url{https://doi.org/10.1090/S0002-9947-03-03366-X}.

\bibitem{Serre}
Serre, J.-P.: \emph{Sur les repr\'esentations modulaires de degr\'e 2 de $\operatorname{Gal}(\overline{\mathbb Q}/\mathbb Q)$}. Duke Math. J. 54(1), 179--230 (1987). \url{https://doi.org/10.1215/S0012-7094-87-05413-5}.

\bibitem{Siksek}
Siksek, S.: \emph{Infinite descent on elliptic curves}. Rocky Mountain J. Math. 25(4), 1501--1538 (1995). \url{https://doi.org/10.1216/rmjm/1181072159}.

\bibitem{SilvermanAdvanced}
Silverman, J.H.: \emph{Advanced Topics in the Arithmetic of Elliptic Curves}. Graduate Texts in Mathematics, vol. 151. Springer, New York (1994). \url{https://doi.org/10.1007/978-1-4612-0851-8}.

\bibitem{Silverman}
Silverman, J.H.: \emph{The Arithmetic of Elliptic Curves}. 2nd edn. Graduate Texts in Mathematics, vol. 106. Springer, Dordrecht (2009). \url{https://doi.org/10.1007/978-0-387-09494-6}.

\bibitem{Simon}
Simon, D.: \emph{Computing the rank of elliptic curves over number fields}. LMS J. Comput. Math. 5, 7--17 (2002). \url{https://doi.org/10.1112/S1461157000000668}.

\bibitem{Sutherland}
Sutherland, A.V.: \emph{Classical modular polynomials}. Electronic coefficient data. Accessed 13 November 2026. \url{https://math.mit.edu/~drew/ClassicalModPolys.html}.

\bibitem{Tate}
Tate, J.: \emph{Algorithm for determining the type of a singular fiber in an elliptic pencil}. In: Birch, B.J., Kuyk, W. (eds.) Modular Functions of One Variable IV, Lecture Notes in Mathematics, vol. 476, pp. 33--52. Springer, Berlin (1975). \url{https://doi.org/10.1007/BFb0097582}.

\bibitem{TaylorWiles}
Taylor, R., Wiles, A.: \emph{Ring-theoretic properties of certain Hecke algebras}. Ann. of Math. (2) 141(3), 553--572 (1995). \url{https://doi.org/10.2307/2118560}.

\bibitem{Velu}
V\'elu, J.: \emph{Isog\'enies entre courbes elliptiques}. C. R. Acad. Sci. Paris Sér. A-B 273, A238--A241 (1971). \url{https://gallica.bnf.fr/ark:/12148/bpt6k56191248/f52.item}.

\bibitem{Wiles}
Wiles, A.: \emph{Modular elliptic curves and Fermat's Last Theorem}. Ann. of Math. (2) 141(3), 443--551 (1995). \url{https://doi.org/10.2307/2118559}.
\end{thebibliography}

% END appendices.tex

\end{document}
